\section*{Ejercicio 5}
\noindent Halle y sombree el área, usando integrales en curvas polares, de la región encerrada por la curva \(r = -12\cos(\theta)\) (circunferencia).

\begin{gather*}
  A / 2 = \frac{1}{2} \int_{0}^{\pi} (-12\cos(\theta))^2 \, d\theta %
    = 72 \int_{0}^{\pi} \cos^2(\theta) \, d\theta %
    = 36 \int_{0}^{\pi} 1 + \cos(2\theta) \, d\theta \\
  A / 2 = 36 \left(\theta + \frac{1 - \sin(2\theta)}{2}\right) \bigg|_{0}^{\pi} %
    = 36 \left(\pi + \frac{1 - \sin(2\pi)}{2}\right) - 36 \left(0 + \frac{1 - \sin(0)}{2}\right) \\
  A / 2 = 36 \left(\pi + \frac{1}{2}\right) - 36 \left(\frac{1}{2}\right) %
    = 36\pi \\
  A = 72\pi
\end{gather*}

\noindent El área de la región encerrada por la curva \(r = -12\cos(\theta)\) es \(72\pi\).

\vspace{4cm}